\scanpage{23}

\deleted{$E$ topos

$C$ petite Catégorie

alors $\uline{\mathrm{Hom}}\,(C^\circ,E)=\pi(C,E)$ un topos

Si on a des morphismes $C\xrightarrow{f} C'$\quad $g:E\longrightarrow E'$ d'où $g^*:E'\longrightarrow E$, on construit un foncteur $\pi(\unsure{h|f},g):\pi(C,E)\longrightarrow\pi(C',E')$
\[
 \pi(f,g)^*=(u'\longmapsto g^*\circ u'\circ f^\circ),\quad
 \pi(f,g)_*=(u'\longmapsto g_*f_*^{(E)})
\]
\margin{vérifier l'expression donnée de $\pi(f,g)_*$ \ldots}
% slanted
qui est le foncteur image inverse associé à un morphisme de topos
\[
 \pi(f,g):\qquad \pi(C,E)\longrightarrow\pi(\unsure{E|C'},E')
\]
de sorte que $\pi(C,E)$ dépend de façon covariante de $(C,E)$ [on a un 2-foncteur $(\mathrm{Cat})\times(\unsure{\mathrm{Top}|\mathscr{T}op})\longrightarrow(\mathrm{Top})$]

On va définir un morphisme \inserted{de topos} fonctoriel en $C,E$
\diagram{\begin{tikzcd}[ampersand replacement=\&]
\text{\textcircled{1}}\& \pi(C,E)\arrow[d,"{\gamma=\gamma(C,E)}"']\\
\&\widehat C\times_{\mathrm{Top}} E
\end{tikzcd}}
en définissant les composantes
\[
 \begin{aligned}
 \gamma_1=\gamma_1(C,E)&:\quad\pi(\,C,E)\longrightarrow\widehat C\\
 \gamma_2=\gamma_2(C,E)&:\quad\pi(C,E)\longrightarrow E
 \end{aligned}
\]
par
\[
 \gamma_1^*:\quad\underset{\substack{\parallel\\\pi(C,P)}}{\widehat C}
 \longrightarrow\pi(C,E)\qquad\text{défini par}\quad F\longmapsto[X\longmapsto F(X)_E]
\]
\inserted{topos ponctuel}\qquad i.e. $\gamma_1=\pi(\mathrm{id}_C,p_E)$, où $p_E:E\longrightarrow P$ \inserted{``unique'' morphisme de topos}
\[
 \gamma_2^*:\quad\underset{\substack{\unsure{\wr|\parallel|\simeq}\\\pi(C_0,\unsure{P|p|E})}}{E}
 \longrightarrow\pi(C,E)\qquad x\longmapsto(i\longmapsto x)
\]
$\gamma_2^*(x)$ est le foncteur constant de valeur $x$

\inserted{catégorie ponctuelle}\qquad i.e. $\gamma_{\unsure{1|2}}=\pi(p_C,\mathrm{id}_E)$,\quad $p_C:C\longrightarrow C_0$ \inserted{unique morphisme de topos}

\uline{Question} \quad $\gamma(C,E)$ est-il tjrs une équivalence ?

a) \deleted{\unsure{Exemples|Exemple} Supposons} On note que $\pi(C,E)$ transforme sommes en un des deux arguments en sommes, et de même pour $\widehat C\times E$ ($C\longrightarrow\widehat C$ transformant sommes en sommes). Donc on est ramené à traiter le cas $C$ connexe.

b) $C$ un groupoïde --- qu'on peut supposer connexe. Alors $\widehat C\simeq B_G$, $G$ un groupe, $\pi(C,E)\simeq B_{G_E}$ (catégorie des objets de $E$ avec $G$ opérant), ok

c) $E=\widehat D$, alors $\pi(C,E)\simeq\widehat{C\times D}$, et $\gamma(C,E)$ est le morphe. \deleted{fonctoriel}\inserted{de topos} évident
\[
 \widehat{C\times{D}}\longrightarrow\widehat C\times\widehat{D}
\]
}
% archive sheet number: 22
